% --- LaTeX Homework Template - S. Venkatraman ---

% --- Set document class and font size ---

\documentclass[letterpaper, 11pt]{article}

% --- Package imports ---

\usepackage{
  amsmath, amsthm, amssymb, mathtools, dsfont,	  % Math typesetting
  graphicx, wrapfig, subfig, float,                  % Figures and graphics formatting
  listings, color, inconsolata, pythonhighlight,     % Code formatting
  fancyhdr, sectsty, hyperref, enumerate, enumitem } % Headers/footers, section fonts, links, lists

% --- Page layout settings ---
\usepackage{bm}
% Set page margins
\usepackage[left=1.35in, right=1.35in, bottom=1in, top=1.1in, headsep=0.2in]{geometry}

% Anchor footnotes to the bottom of the page
\usepackage[bottom]{footmisc}

% Set line spacing
\renewcommand{\baselinestretch}{1.2}

% Set spacing between paragraphs
\setlength{\parskip}{1.5mm}

% Allow multi-line equations to break onto the next page
\allowdisplaybreaks

% Enumerated lists: make numbers flush left, with parentheses around them
\setlist[enumerate]{wide=0pt, leftmargin=21pt, labelwidth=0pt, align=left}
\setenumerate[1]{label={(\arabic*)}}

% --- Page formatting settings ---

% Set link colors for labeled items (blue) and citations (red)
\hypersetup{colorlinks=true, linkcolor=blue, citecolor=red}

% Make reference section title font smaller
\renewcommand{\refname}{\large\bf{References}}

% --- Settings for printing computer code ---

% Define colors for green text (comments), grey text (line numbers),
% and green frame around code
\definecolor{greenText}{rgb}{0.5, 0.7, 0.5}
\definecolor{greyText}{rgb}{0.5, 0.5, 0.5}
\definecolor{codeFrame}{rgb}{0.5, 0.7, 0.5}

% Define code settings
\lstdefinestyle{code} {
  frame=single, rulecolor=\color{codeFrame},            % Include a green frame around the code
  numbers=left,                                         % Include line numbers
  numbersep=8pt,                                        % Add space between line numbers and frame
  numberstyle=\tiny\color{greyText},                    % Line number font size (tiny) and color (grey)
  commentstyle=\color{greenText},                       % Put comments in green text
  basicstyle=\linespread{1.1}\ttfamily\footnotesize,    % Set code line spacing
  keywordstyle=\ttfamily\footnotesize,                  % No special formatting for keywords
  showstringspaces=false,                               % No marks for spaces
  xleftmargin=1.95em,                                   % Align code frame with main text
  framexleftmargin=1.6em,                               % Extend frame left margin to include line numbers
  breaklines=true,                                      % Wrap long lines of code
  postbreak=\mbox{\textcolor{greenText}{$\hookrightarrow$}\space} % Mark wrapped lines with an arrow
}

% Set all code listings to be styled with the above settings
\lstset{style=code}

% --- Math/Statistics commands ---

% Add a reference number to a single line of a multi-line equation
% Usage: "\numberthis\label{labelNameHere}" in an align or gather environment
\newcommand\numberthis{\addtocounter{equation}{1}\tag{\theequation}}

% Shortcut for bold text in math mode, e.g. $\b{X}$
\let\b\mathbf

% Shortcut for bold Greek letters, e.g. $\bg{\beta}$
\let\bg\boldsymbol

% Shortcut for calligraphic script, e.g. %\mc{M}$
\let\mc\mathcal

% \mathscr{(letter here)} is sometimes used to denote vector spaces
\usepackage[mathscr]{euscript}

% Convergence: right arrow with optional text on top
% E.g. $\converge[w]$ for weak convergence
\newcommand{\converge}[1][]{\xrightarrow{#1}}

% Normal distribution: arguments are the mean and variance
% E.g. $\normal{\mu}{\sigma}$
\newcommand{\normal}[2]{\mathcal{N}\left(#1,#2\right)}

% Uniform distribution: arguments are the left and right endpoints
% E.g. $\unif{0}{1}$
\newcommand{\unif}[2]{\text{Uniform}(#1,#2)}

% Independent and identically distributed random variables
% E.g. $ X_1,...,X_n \iid \normal{0}{1}$
\newcommand{\iid}{\stackrel{\smash{\text{iid}}}{\sim}}

% Equality: equals sign with optional text on top
% E.g. $X \equals[d] Y$ for equality in distribution
\newcommand{\equals}[1][]{\stackrel{\smash{#1}}{=}}

% Math mode symbols for common sets and spaces. Example usage: $\R$
\newcommand{\R}{\mathbb{R}}   % Real numbers
\newcommand{\C}{\mathbb{C}}   % Complex numbers
\newcommand{\Q}{\mathbb{Q}}   % Rational numbers
\newcommand{\Z}{\mathbb{Z}}   % Integers
\newcommand{\N}{\mathbb{N}}   % Natural numbers
\newcommand{\F}{\mathcal{F}}  % Calligraphic F for a sigma algebra
\newcommand{\El}{\mathcal{L}} % Calligraphic L, e.g. for L^p spaces

% Math mode symbols for probability
\newcommand{\pr}{\mathbb{P}}    % Probability measure
\newcommand{\E}{\mathbb{E}}     % Expectation, e.g. $\E(X)$
\newcommand{\var}{\text{Var}}   % Variance, e.g. $\var(X)$
\newcommand{\cov}{\text{Cov}}   % Covariance, e.g. $\cov(X,Y)$
\newcommand{\corr}{\text{Corr}} % Correlation, e.g. $\corr(X,Y)$
\newcommand{\B}{\mathcal{B}}    % Borel sigma-algebra

% Other miscellaneous symbols
\newcommand{\tth}{\text{th}}	% Non-italicized 'th', e.g. $n^\tth$
\newcommand{\Oh}{\mathcal{O}}	% Big-O notation, e.g. $\O(n)$
\newcommand{\1}{\mathds{1}}	% Indicator function, e.g. $\1_A$

% Additional commands for math mode
\DeclareMathOperator*{\argmax}{argmax}    % Argmax, e.g. $\argmax_{x\in[0,1]} f(x)$
\DeclareMathOperator*{\argmin}{argmin}    % Argmin, e.g. $\argmin_{x\in[0,1]} f(x)$
\DeclareMathOperator*{\spann}{Span}       % Span, e.g. $\spann\{X_1,...,X_n\}$
\DeclareMathOperator*{\bias}{Bias}        % Bias, e.g. $\bias(\hat\theta)$
\DeclareMathOperator*{\ran}{ran}          % Range of an operator, e.g. $\ran(T) 
\DeclareMathOperator*{\dv}{d\!}           % Non-italicized 'with respect to', e.g. $\int f(x) \dv x$
\DeclareMathOperator*{\diag}{diag}        % Diagonal of a matrix, e.g. $\diag(M)$
\DeclareMathOperator*{\trace}{trace}      % Trace of a matrix, e.g. $\trace(M)$

% Numbered theorem, lemma, etc. settings - e.g., a definition, lemma, and theorem appearing in that 
% order in Section 2 will be numbered Definition 2.1, Lemma 2.2, Theorem 2.3. 
% Example usage: \begin{theorem}[Name of theorem] Theorem statement \end{theorem}
\theoremstyle{definition}
\newtheorem{theorem}{Theorem}[section]
\newtheorem{proposition}[theorem]{Proposition}
\newtheorem{lemma}[theorem]{Lemma}
\newtheorem{corollary}[theorem]{Corollary}
\newtheorem{definition}[theorem]{Definition}
\newtheorem{example}[theorem]{Example}
\newtheorem{remark}[theorem]{Remark}

% Un-numbered theorem, lemma, etc. settings
% Example usage: \begin{lemma*}[Name of lemma] Lemma statement \end{lemma*}
\newtheorem*{theorem*}{Theorem}
\newtheorem*{proposition*}{Proposition}
\newtheorem*{lemma*}{Lemma}
\newtheorem*{corollary*}{Corollary}
\newtheorem*{definition*}{Definition}
\newtheorem*{example*}{Example}
\newtheorem*{remark*}{Remark}
\newtheorem*{claim}{Claim}

\newcommand{\Pp}{\mathbb{P}}
\newcommand{\Dkl}{\mathrm{D}}

% --- Left/right header text (to appear on every page) ---

% Include a line underneath the header, no footer line
\pagestyle{fancy}
\renewcommand{\footrulewidth}{0pt}
\renewcommand{\headrulewidth}{0.4pt}

% Left header text: course name/assignment number
\lhead{Foundations of Machine Learning-- Homework 1}

% Right header text: your name
\rhead{Yixia Yu (yy5091@nyu.edu)}

% --- Document starts here ---

\begin{document}
\subsection*{Problem A. Concentration inequality}
In this problem, we will derive a concentration inequality that is finer than Hoeffding's inequality.

Let $X_1,\ldots,X_m$ be independent random variables drawn according to some distribution $\mathcal{D}$ with mean $p$ and support included in $[0,1]$. Let $q\in[0,1]$ and define $\hat p$ by $\hat p=\frac{1}{m}\sum_{i=1}^m X_i$.

\begin{enumerate}[leftmargin=1.2em,label=\arabic*.]
\item Use Markov's inequality to prove that for any $t>0$,
\[
\Pp[\hat p\ge q]\le e^{-tm q}\prod_{i=1}^m \E\!\left[e^{t X_i}\right].
\]

\item Show that
\[
\Pp[\hat p\ge q]\le \bigl[e^{-t q}\,\bigl(1-p+e^{t}p\bigr)\bigr]^m
\]
\emph{(Hint: you can show that $\forall x\in[0,1]$, $e^{t x}\le 1-x+e^{t}x$).}

\item Choose $t>0$ to minimize the right-hand side and prove the following inequality:
\[
\Pp[\hat p\ge q]\le e^{-m\,\Dkl(q\|p)},
\]
where $\Dkl(q\|p)= q\log\frac{q}{p}+(1-q)\log\frac{1-q}{1-p}$ is the binary relative entropy of $p$ and $q$.

\item Show that this can be used to derive a bound that is finer than Hoeffding's inequality by choosing $q=p+\varepsilon$, for any $\varepsilon>0$ \emph{(Hint: you can use Pinsker's inequality, which in this context gives $\Dkl(q\|p)\ge 2|p-q|^2$)}.
\end{enumerate}

(1)Recall that if \(Y\) is a nonnegative random variable and \(a > 0\), then Markov's inequality states:
   \[
   \mathbb{P}[Y \ge a] \le \frac{\mathbb{E}[Y]}{a}.
   \]
   Here, let
   \[
   Y = e^{tm\hat{p}} \quad \text{and} \quad a = e^{tmq}.
   \]
   Then,
   \[
   \mathbb{P}\left(e^{tm\hat{p}} \ge e^{tmq}\right) \le \frac{\mathbb{E}\left[e^{tm\hat{p}}\right]}{e^{tmq}}.
   \]
   Since the event \(\{e^{tm\hat{p}} \ge e^{tmq}\}\) is equivalent to \(\{\hat{p} \ge q\}\), we have:
   \[
   \mathbb{P}\left(\hat{p} \ge q\right) \le e^{-tmq} \, \mathbb{E}\left[e^{tm\hat{p}}\right].
   \]
Notice that
   \[
   \hat{p} = \frac{1}{m} \sum_{i=1}^m X_i \quad \Longrightarrow \quad e^{tm\hat{p}} = e^{t \sum_{i=1}^m X_i}.
   \]
   Since the \(X_i\) are independent, the expectation of the exponential of the sum factors as a product:
   \[
   \mathbb{E} \left[e^{t\sum_{i=1}^m X_i}\right] = \prod_{i=1}^m \mathbb{E} \left[e^{tX_i}\right].
   \]
   Thus, we obtain:
   \[
   \mathbb{P}\left(\hat{p} \ge q\right) \le e^{-tmq} \prod_{i=1}^m \mathbb{E} \left[e^{tX_i}\right].
   \]
(2)Define the function
\[
f(x) = e^{tx}.
\]
We compute its second derivative:
\[
f''(x) = t^2 e^{tx}.
\]
Since \(t > 0\) and \(e^{tx} > 0\) for all \(x\), we have \(f''(x) > 0\) for all \(x\). This shows that \(f(x)\) is a \emph{convex} function.

For any \(x\in [0,1]\), we can express \(x\) as a convex combination of 0 and 1:
\[
x = (1-x)\cdot 0 + x\cdot 1.
\]
Using the property of convex functions, we have:
\[
f\Bigl((1-x)\cdot 0 + x\cdot 1\Bigr) \le (1-x)f(0) + x f(1).
\]
the inequality becomes
\[
e^{tx} \le (1-x)\cdot 1 + x\cdot e^t = 1-x+e^t x.
\]
This completes the proof of hint.
\\Taking expectation with respect to \(X_i\) yields
\[
\mathbb{E}\left[e^{tX_i}\right] \le \mathbb{E}\Bigl[1-X_i+e^t X_i\Bigr].
\]
Because expectation is linear, we have
\[
\mathbb{E}\Bigl[1-X_i+e^t X_i\Bigr] = 1-\mathbb{E}[X_i]+e^t\,\mathbb{E}[X_i] = 1-p+e^t p.
\]
Since the same bound holds for each \(X_i\), we get
\[
\prod_{i=1}^m \mathbb{E}\left[e^{tX_i}\right] \le \Bigl(1-p+e^t p\Bigr)^m.
\]
Substituting this into our earlier bound gives
\[
\mathbb{P}\left[\hat{p}\ge q\right] \le e^{-tmq}\Bigl(1-p+e^t p\Bigr)^m.
\]
Since \(e^{-tmq} = \Bigl(e^{-tq}\Bigr)^m\), we can write:
\[
\mathbb{P}\left[\hat{p}\ge q\right] \le \Bigl[e^{-tq}\Bigl(1-p+e^t p\Bigr)\Bigr]^m.
\]
This completes the proof.
\\(3)
\\
Let 
\[
f(t)= e^{-tq}\Bigl(1-p+p\,e^t\Bigr).
\]
We wish to choose \(t>0\) so that \(f(t)\) is minimized.
\\Define
\[
\phi(t)=\ln f(t)= -tq + \ln\Bigl(1-p+p\,e^t\Bigr).
\]
Differentiating with respect to \(t\) gives
\[
\phi'(t) = -q + \frac{p\,e^t}{\,1-p+p\,e^t}\,.
\]
Setting \(\phi'(t)=0\) yields
\[
\frac{p\,e^t}{\,1-p+p\,e^t} = q.
\]
Solving for \(e^t\):
\[
p\,e^t = q\Bigl(1-p+p\,e^t\Bigr) \quad \Longrightarrow \quad p\,e^t = q(1-p) + q\,p\,e^t\,.
\]
Rearrange the terms to isolate \(e^t\):
\[
p\,e^t (1-q) = q(1-p) \quad \Longrightarrow \quad e^t = \frac{q(1-p)}{p(1-q)}.
\]
Thus, the minimizing value is
\[
t^* = \ln\!\left(\frac{q(1-p)}{p(1-q)}\right).
\]

Substitute \(t=t^*\) into the bound:
\[
f(t^*) = e^{-q\,t^*}\Bigl(1-p+p\,e^{t^*}\Bigr).
\]
Since
\[
e^{t^*} = \frac{q(1-p)}{p(1-q)},
\]
we have:
\[
e^{-q\,t^*} = \left(\frac{q(1-p)}{p(1-q)}\right)^{-q},
\]
and
\[
1-p+p\,e^{t^*}= 1-p + p\,\frac{q(1-p)}{p(1-q)}
= 1-p + \frac{q(1-p)}{1-q}
= (1-p)\left(1 + \frac{q}{1-q}\right)
= \frac{1-p}{1-q}\,.
\]
Thus,
\[
f(t^*) = \left(\frac{q(1-p)}{p(1-q)}\right)^{-q}\frac{1-p}{1-q}.
\]

Rewrite the expression 
and we obtain
\[
f(t^*) = \left(\frac{p}{q}\right)^q (1-p)^{1-q}(1-q)^{-(1-q)}
= \left(\frac{p}{q}\right)^q \left(\frac{1-p}{1-q}\right)^{1-q}\,.
\]
Observing that
\[
e^{-\Dkl(q\|p)} = \exp\Bigl[-q\log\frac{q}{p}-(1-q)\log\frac{1-q}{1-p}\Bigr]
= \left(\frac{p}{q}\right)^q \left(\frac{1-p}{1-q}\right)^{1-q}\,,
\]
we conclude that
\[
f(t^*) = e^{-\Dkl(q\|p)}.
\]
Thus, the optimal bound becomes
\[
\mathbb{P}\Bigl[\hat{p}\ge q\Bigr] \le \Bigl[f(t^*)\Bigr]^m
= e^{-m\,\Dkl(q\|p)}\,.
\]
\\(4)
choose 
\[
q = p + \varepsilon \quad \text{for any } \varepsilon > 0.
\]
Then the bound becomes
\[
\mathbb{P}\left[\hat{p} \ge p+\varepsilon\right] \le \exp\Bigl(-m\,\Dkl(p+\varepsilon\|p)\Bigr).
\]

\emph{Pinsker's inequality}, tells us that in this context
\[
\Dkl(p+\varepsilon \| p) \ge 2\varepsilon^2.
\]
Substituting this into our bound, we obtain
\[
\mathbb{P}\left[\hat{p} \ge p+\varepsilon\right] \le e^{-2m\,\varepsilon^2}.
\]

This bound is at least as strong (and under many circumstances finer) than Hoeffding's inequality.  because:
\begin{itemize}
  \item The binary relative entropy \(\Dkl(p+\varepsilon\|p)\) depends on the specific values of \(p\) and \(p+\varepsilon\), thereby capturing the true large-deviation rate. In contrast, Hoeffding's inequality provides a uniform \emph{worst-case} bound over all distributions with support in \([0,1]\).
  \item Pinsker's inequality gives the lower bound \(\Dkl(p+\varepsilon\|p) \ge 2\varepsilon^2\), but \(\Dkl(p+\varepsilon\|p)\) is typically \emph{greater} than \(2\varepsilon^2\) unless \(p=1/2\) (or in the limit as \(\varepsilon\to 0\)). For example, when \(p\) is near 0 or 1, the divergence \(\Dkl(p+\varepsilon\|p)\) can be significantly larger than \(2\varepsilon^2\), which implies a much faster exponential decay for the tail probability.
\end{itemize}

\subsection*{Problem B. PAC-Learning}
\begin{enumerate}[leftmargin=1.2em,label=\arabic*.]
\item \textbf{Concentric spheres.} Let $\mathcal{X}=\mathbb{R}^3$ and consider the set of concepts of the form
\[
c=\{(x,y,z): x^2+y^2+z^2\le r^2\}\quad\text{for some } r>0.
\]
Show that this class can be $(\varepsilon,\delta)$-PAC-learned from training data of size
\[
m \;\ge\; \frac{1}{\varepsilon}\,\log\!\left(\frac{1}{\delta}\right).
\]
You should carefully justify all steps of your proof. (For this problem, you should seek a proof similar to one presented in class for rectangles in the plane and not leverage a more powerful or more general theorem.)
\end{enumerate}
We consider the concept class
\[
\mathcal{C} = \{ c_r : c_r = \{(x,y,z) \in \mathbb{R}^3 : x^2 + y^2 + z^2 \le r^2\}, \; r > 0\},
\]
and assume the target concept is
\[
c = c_{r_0} = \{(x,y,z) \in \mathbb{R}^3 : x^2 + y^2 + z^2 \le r_0^2\},
\]
for some unknown \(r_0 > 0\). The distribution \(D\) over \(\mathbb{R}^3\) is arbitrary.

Then, we can he this Algorithm: Given \(m\) i.i.d. training examples
\[
S = \{(x_1, y_1), (x_2, y_2), \ldots, (x_m, y_m)\},
\]
where \(x_i \sim D\) and
\[
y_i = \mathbf{1}_{x_i \in c},
\]
the algorithm outputs the hypothesis
\[
h = \{ x \in \mathbb{R}^3 : \|x\| \le r_+\},
\]
where
\[
r_+ = \max\{\, \|x_i\| : y_i = 1 \,\}
\]
is the maximum distance from the origin among all positive examples.

The correctness is that since for any positive example we have \(\|x_i\| \le r_0\), it follows that
\[
r_+ \le r_0,
\]
so that
\[
h \subseteq c.
\]
Thus, the hypothesis \(h\) never labels a negative example as positive and the only possible error is failing to cover some positive examples. In particular, the error is
\[
\operatorname{error}_D(h) = \Pr_{x \sim D}\big[x \in c \setminus h\big].
\]

\textbf{PAC Guarantee}:\\
Fix \(\varepsilon,\delta \in (0,1)\). We wish to show that if
\[
m \ge \frac{1}{\varepsilon}\ln\!\left(\frac{1}{\delta}\right),
\]
then
\[
\Pr\big[\operatorname{error}_D(h) > \varepsilon\big] \le \delta.
\]

The proof proceeds similarly to the one for axis-aligned rectangles in the plane:

Since \(h \subseteq c\), the error comes solely from points in the annulus (the outer shell) of \(c\). Choose a radius \(r'\) with \(0 \le r' < r_0\) so that the outer region
\[
A = \{x \in c: r' < \|x\| \le r_0\}
\]
has probability
\[
\Pr_D[A] \ge \varepsilon.
\]
(If such an \(A\) does not exist—namely, if \(\Pr_D[c] \le \varepsilon\)—then for any hypothesis \(h \subseteq c\) the error satisfies \(\operatorname{error}_D(h) \le \Pr_D[c] \le \varepsilon\), and the claim holds trivially.)

Notice that if at least one positive training example lies in \(A\), then the maximum radius \(r_+\) will be at least \(r'\), and therefore
\[
h = \{x: \|x\| \le r_+\} \supseteq \{x: \|x\| \le r'\}.
\]
Thus,
\[
c \setminus h \subseteq c \setminus \{x: \|x\| \le r'\} = A,
\]
and consequently,
\[
\operatorname{error}_D(h) \le \Pr_D[A] \le \varepsilon.
\]
In other words, if a positive training example falls in \(A\), then the error is at most \(\varepsilon\).

If \(\operatorname{error}_D(h) > \varepsilon\), then no positive training example in \(S\) fell in \(A\). For each training example, the probability (unconditionally) that it lands in \(A\) is at least \(\varepsilon\). Therefore, the probability that a single training example does \emph{not} fall in \(A\) is at most \(1-\varepsilon\). Since the examples are independent, the probability that none of the \(m\) examples falls in \(A\) is bounded by
\[
(1-\varepsilon)^{\,m} \le e^{-m\varepsilon}.
\]

To guarantee that
\[
\Pr\big[\operatorname{error}_D(h) > \varepsilon\big] \le e^{-m\varepsilon} \le \delta,
\]
it suffices to choose \(m\) so that
\[
e^{-m\varepsilon} \le \delta.
\]
Taking logarithms we have
\[
-m\varepsilon \le \ln \delta,
\]
or equivalently,
\[
m \ge \frac{1}{\varepsilon}\ln\!\left(\frac{1}{\delta}\right).
\]

This completes the proof
\hfill \(\square\)
\subsection*{Problem C. Rademacher complexity}
\begin{enumerate}[leftmargin=1.2em,label=\arabic*.]
\item Let $\mathcal{H}_1$ and $\mathcal{H}_2$ be two families of functions mapping $\mathcal{X}$ to $[0,1]$ and let $S$ be a sample of size $m$.
  \begin{enumerate}[label=(\alph*)]
  \item Show that for any sample $S$,
  \[
  \widehat{\mathfrak{R}}_{S}(\mathcal{H}_1+\mathcal{H}_2)
  =\widehat{\mathfrak{R}}_{S}(\mathcal{H}_1)+\widehat{\mathfrak{R}}_{S}(\mathcal{H}_2).
  \]
  \item Show that the empirical Rademacher complexity of $\mathcal{H}$ for any sample $S$ of size $m$ can be bounded as follows:
  \[
  \widehat{\mathfrak{R}}_{S}(\mathcal{H}_1\mathcal{H}_2)
  \le \frac{3}{2}\Bigl[\widehat{\mathfrak{R}}_{S}(\mathcal{H}_1)+\widehat{\mathfrak{R}}_{S}(\mathcal{H}_2)\Bigr].
  \]
  \textit{Hint:} use the identity $x_1 x_2=\tfrac14\bigl[(x_1+x_2)^2-(x_1-x_2)^2\bigr]$ valid for all $x_1,x_2\in\mathbb{R}$ and Talagrand's contraction lemma.
  \end{enumerate}

\item Let $\mathcal{E}$ be the family of functions defined over $\mathcal{X}$ by $\mathcal{E}=\{x\mapsto \mathbb{1}_{x=z}\colon z\in\mathcal{X}\}$.
  \begin{enumerate}[label=(\alph*)]
  \item Let $S$ be a sample of size $m$ with distinct points. Prove the inequality $\widehat{\mathfrak{R}}_{S}(\mathcal{E})\le \tfrac{1}{m}$.
  \item Let $S$ be a sample of size $m$ with a single point repeated $m$ times. Prove the equality
  \[
  \widehat{\mathfrak{R}}_{S}(\mathcal{E})=\frac{1}{2m}\,\mathbb{E}_{\boldsymbol{\sigma}}\bigl[\left| \sum_{i=1}^{m}\sigma_i \right|\bigr].
  \]
  Show that $\widehat{\mathfrak{R}}_{S}(\mathcal{E})\le \tfrac{1}{2\sqrt{m}}$.
  \item  1[bonus question] Let $n^\ast$ denote the largest number of repetitions of any point in the sample. Give an upper bound on $\widehat{\mathfrak{R}}_{S}(\mathcal{E})$ in terms of $n^\ast$.
  \end{enumerate}
\end{enumerate}

(1)
For the class \(\mathcal{H}_1+\mathcal{H}_2\), we have
\[
\begin{aligned}
\widehat{\mathfrak{R}}_S(\mathcal{H}_1+\mathcal{H}_2) 
&=\frac{1}{m}\mathbb{E}_\sigma \left[\sup_{h_1\in\mathcal{H}_1,\,h_2\in\mathcal{H}_2} \sum_{i=1}^m \sigma_i \Bigl(h_1(x_i)+h_2(x_i)\Bigr)\right] \\
&=\frac{1}{m}\mathbb{E}_\sigma \Biggl[
\sup_{h_1\in\mathcal{H}_1}\sum_{i=1}^m \sigma_i h_1(x_i)
+\sup_{h_2\in\mathcal{H}_2}\sum_{i=1}^m \sigma_i h_2(x_i)
\Biggr] \\
&=\widehat{\mathfrak{R}}_S(\mathcal{H}_1)+\widehat{\mathfrak{R}}_S(\mathcal{H}_2).
\end{aligned}
\]
 This completes the proof for part (a).

Let us now consider the class \(\mathcal{H}_1\mathcal{H}_2=\{\,h_1\cdot h_2: h_1\in\mathcal{H}_1,\,h_2\in\mathcal{H}_2\,\}\). By definition,
\[
\widehat{\mathfrak{R}}_S(\mathcal{H}_1\mathcal{H}_2)
=\frac{1}{m}\mathbb{E}_\sigma \left[\sup_{h_1\in\mathcal{H}_1,\, h_2\in\mathcal{H}_2} \sum_{i=1}^m \sigma_i \, h_1(x_i)h_2(x_i)\right].
\]
For any real numbers \(x_1,x_2\) the following identity holds:
\[
x_1x_2=\frac{1}{4}\Bigl[(x_1+x_2)^2-(x_1-x_2)^2\Bigr].
\]
Hence, replacing \(x_1=r_i:=h_1(x_i)\) and \(x_2=s_i:=h_2(x_i)\) (noting that since \(h_1(x_i),h_2(x_i)\in[0,1]\) the sums \(h_1(x_i)+h_2(x_i)\in[0,2]\) and differences \(\;h_1(x_i)-h_2(x_i)\in[-1,1]\)) we obtain
\[
h_1(x_i)h_2(x_i)=\frac{1}{4}\Bigl[(h_1(x_i)+h_2(x_i))^2-(h_1(x_i)-h_2(x_i))^2\Bigr].
\]
Thus, for any fixed choice of Rademacher signs,
\[
\sum_{i=1}^m\sigma_i\, h_1(x_i)h_2(x_i)
=\frac{1}{4}\sum_{i=1}^m\sigma_i\Bigl[(h_1(x_i)+h_2(x_i))^2-(h_1(x_i)-h_2(x_i))^2\Bigr].
\]
Taking the supremum over \(h_1\) and \(h_2\) and then the expectation yields
\[
\widehat{\mathfrak{R}}_S(\mathcal{H}_1\mathcal{H}_2)
\le \frac{1}{4}\left[
\widehat{\mathfrak{R}}_S\Bigl(\{(h_1+h_2)^2\}\Bigr)
+\widehat{\mathfrak{R}}_S\Bigl(\{(h_1-h_2)^2\}\Bigr)
\right].
\]

We now use Talagrand's contraction lemma. Recall that if \(\phi\colon\mathbb{R}\to\mathbb{R}\) is \(L\)-Lipschitz with \(\phi(0)=0\), then
\[
\widehat{\mathfrak{R}}_S\Bigl(\phi\circ\mathcal{F}\Bigr)
\le L\,\widehat{\mathfrak{R}}_S(\mathcal{F}).
\]

\medskip

\underline{For the Class \(\{(h_1+h_2)^2\}\):}  
Since \(h_1(x)+h_2(x)\in[0,2]\), the function \(\phi(x)=x^2\) is \(4\)-Lipschitz on \([0,2]\) (because for \(x,y\in[0,2]\),
\[
|x^2-y^2|=|x-y||x+y|\le 4|x-y|).
\]
Moreover, \(\phi(0)=0\). Thus, by the contraction lemma,
\[
\widehat{\mathfrak{R}}_S\Bigl(\{(h_1+h_2)^2\}\Bigr)
\le 4\,\widehat{\mathfrak{R}}_S(\{h_1+h_2\}).
\]
Using part (a), we have
\[
\widehat{\mathfrak{R}}_S(\{h_1+h_2\})=\widehat{\mathfrak{R}}_S(\mathcal{H}_1)
+\widehat{\mathfrak{R}}_S(\mathcal{H}_2).
\]
Thus,
\[
\widehat{\mathfrak{R}}_S\Bigl(\{(h_1+h_2)^2\}\Bigr)
\le 4\,\Bigl[\widehat{\mathfrak{R}}_S(\mathcal{H}_1)
+\widehat{\mathfrak{R}}_S(\mathcal{H}_2)\Bigr].
\]

\medskip

\underline{For the Class \(\{(h_1-h_2)^2\}\):}  
Similarly, \(h_1(x)-h_2(x)\in[-1,1]\) and on \([-1,1]\) the function \(\psi(x)=x^2\) is \(2\)-Lipschitz (since \(|x^2-y^2|\le 2|x-y|\) for all \(x,y\in[-1,1]\) and \(\psi(0)=0\)). Hence,
\[
\widehat{\mathfrak{R}}_S\Bigl(\{(h_1-h_2)^2\}\Bigr)
\le 2\,\widehat{\mathfrak{R}}_S(\{h_1-h_2\})
= 2\,\Bigl[\widehat{\mathfrak{R}}_S(\mathcal{H}_1)
+\widehat{\mathfrak{R}}_S(\mathcal{H}_2)\Bigr],
\]
where in the last step we again used part (a) (noting that for any pair, \((h_1-h_2)\) belongs to the sum of \(\mathcal{H}_1\) and \(-\mathcal{H}_2\); since the Rademacher complexity is invariant under sign changes, the same bound holds).

\medskip

Plugging these estimates back into our bound for \(\widehat{\mathfrak{R}}_S(\mathcal{H}_1\mathcal{H}_2)\) we obtain
\[
\begin{aligned}
\widehat{\mathfrak{R}}_S(\mathcal{H}_1\mathcal{H}_2)
&\le \frac{1}{4}\left[
4\,\Bigl(\widehat{\mathfrak{R}}_S(\mathcal{H}_1)
+\widehat{\mathfrak{R}}_S(\mathcal{H}_2)\Bigr)
+ 2\,\Bigl(\widehat{\mathfrak{R}}_S(\mathcal{H}_1)
+\widehat{\mathfrak{R}}_S(\mathcal{H}_2)\Bigr)
\right] \\
&=\frac{1}{4}\cdot 6\,\Bigl(\widehat{\mathfrak{R}}_S(\mathcal{H}_1)
+\widehat{\mathfrak{R}}_S(\mathcal{H}_2)\Bigr)
=\frac{3}{2}\,\Bigl(\widehat{\mathfrak{R}}_S(\mathcal{H}_1)
+\widehat{\mathfrak{R}}_S(\mathcal{H}_2)\Bigr).
\end{aligned}
\]
This completes the proof for part (b).\\
(2)
\\\textbf{Proof:}\\[1mm]
For the class \(\mathcal{E}\), each function is of the form
\[
f_z(x)=\mathbf{1}_{\{x=z\}},\quad \text{for some } z\in\mathcal{X}.
\]
Consider a fixed sample \(S=\{x_1,\dots,x_m\}\) where all \(x_i\) are distinct. There are two types of functions in \(\mathcal{E}\):
  
1. If \(z\notin S\), then for all \(i\) we have \(f_z(x_i)=0\) and thus 
\[
\sum_{i=1}^{m}\sigma_i\,f_z(x_i)=0.
\]
  
2. If \(z=x_j\) for some \(j\in\{1,\dots,m\}\), then
\[
f_z(x_i)= \begin{cases} 1, &\text{if }i=j,\\ 0,& \text{otherwise.} \end{cases}
\]
Hence,
\[
\sum_{i=1}^m \sigma_i\,f_z(x_i)=\sigma_j.
\]

Thus, the supremum 
\[
\sup_{f\in \mathcal{E}}\sum_{i=1}^m \sigma_i\,f(x_i)
\]
is achieved by choosing an index \(j\) for which \(\sigma_j=+1\) if at least one exists; otherwise, if all \(\sigma_i=-1\), one may choose a function corresponding to some \(z\notin S\) (which gives sum \(0\)) to obtain a larger value. In summary, we have
\[
\sup_{f\in \mathcal{E}}\sum_{i=1}^m \sigma_i\,f(x_i)
=\max\Bigl\{0,\,\max_{1\le j\le m} \sigma_j\Bigr\}.
\]
Since each \(\sigma_j\) takes values in \(\{-1,+1\}\), the above supremum equals:
\[
\sup_{f\in \mathcal{E}}\sum_{i=1}^m \sigma_i\,f(x_i)=
\begin{cases}
1, &\text{if there exists some } j \text{ with } \sigma_j=+1; \\
0, &\text{if } \sigma_i=-1 \text{ for all } i.
\end{cases}
\]

Therefore, by taking expectation over the random signs, we observe that
\[
\mathbb{E}_{\sigma}\!\left[\sup_{f\in \mathcal{E}}\sum_{i=1}^m \sigma_i\,f(x_i)\right]
=(1-2^{-m})\cdot 1 + 2^{-m}\cdot 0 = 1-2^{-m}.
\]

Finally, by the definition of empirical Rademacher complexity we have
\[
\widehat{\mathfrak{R}}_S(\mathcal{E})
=\frac{1}{m}\Bigl(1-2^{-m}\Bigr)
\le \frac{1}{m},
\]
since \(1-2^{-m}\le 1\) for all \(m\ge 1\).

This completes the proof.
\\(b)
Since every function in \(\mathcal{E}\) is of the form
\[
f_{z}(x) = \mathbf{1}_{\{x=z\}} \quad \text{for some } z\in\mathcal{X},
\]
we examine its action on the sample \(S\).

Because \(S\) contains only the single point \(x\), we have:
- If \(z=x\): then \(f_x(x) = 1\) so that for every \(i\),
  \[
  f_x(x_i) = 1,
  \]
  and therefore
  \[
  \sum_{i=1}^{m} \sigma_i\, f_x(x_i) = \sum_{i=1}^{m} \sigma_i.
  \]
- For any \(z \neq x\): then \(f_z(x) = 0\) so that
  \[
  f_z(x_i) = 0 \quad \text{for all } i,
  \]
  leading to
  \[
  \sum_{i=1}^{m} \sigma_i\, f_z(x_i) = 0.
  \]

Thus, the supremum over \(\mathcal{E}\) is
\[
\sup_{f\in\mathcal{E}} \sum_{i=1}^m \sigma_i\, f(x_i)
=\max\left\{\sum_{i=1}^m \sigma_i,\; 0\right\},
\]
which is the positive part of the sum. We denote it by
\[
\left(\sum_{i=1}^m \sigma_i\right)_+.
\]

Thus, the empirical Rademacher complexity becomes
\[
\widehat{\mathfrak{R}}_S(\mathcal{E}) 
=\frac{1}{m}\,\mathbb{E}_{\boldsymbol{\sigma}}\!\left[\left(\sum_{i=1}^m \sigma_i\right)_+\right].
\]
Because the Rademacher distribution is symmetric, the positive part and the negative part are balanced. In fact, note that
\[
\left|\sum_{i=1}^m \sigma_i\right| 
=\left(\sum_{i=1}^m \sigma_i\right)_+ + \left(-\sum_{i=1}^m \sigma_i\right)_+
=2\left(\sum_{i=1}^m \sigma_i\right)_+,
\]
where the last equality holds by symmetry (i.e. \(\mathbb{E}[\left(\sum \sigma_i\right)_+]=\mathbb{E}\left[\left(-\sum \sigma_i\right)_+\right]\)). Taking expectation, we obtain
\[
\mathbb{E}_{\boldsymbol{\sigma}}\!\left[\left(\sum_{i=1}^m \sigma_i\right)_+\right]
=\frac{1}{2}\,\mathbb{E}_{\boldsymbol{\sigma}}\!\left[\left|\sum_{i=1}^m \sigma_i\right|\right].
\]
Substituting back, we have
\[
\widehat{\mathfrak{R}}_S(\mathcal{E})
=\frac{1}{2m}\,\mathbb{E}_{\boldsymbol{\sigma}}\!\left[\left|\sum_{i=1}^{m}\sigma_i\right|\right].
\]

To bound \(\mathbb{E}\Bigl[\left|\sum_{i=1}^{m}\sigma_i\right|\Bigr]\), we can use the following standard calculation. First, notice that
\[
\mathbb{E}\!\left[\left(\sum_{i=1}^{m}\sigma_i\right)^2\right]
=\sum_{i=1}^{m} \mathbb{E}[\sigma_i^2] 
= m,
\]
since the \(\sigma_i\) are independent with \(\sigma_i^2=1\) almost surely.
Now, by Jensen’s inequality, we have
\[
\mathbb{E}\!\left[\left|\sum_{i=1}^{m}\sigma_i\right|\right]
\le \sqrt{\mathbb{E}\!\left[\Bigl(\sum_{i=1}^{m}\sigma_i\Bigr)^2\right]}
=\sqrt{m}.
\]
Thus,
\[
\widehat{\mathfrak{R}}_S(\mathcal{E})
\le \frac{1}{2m}\,\sqrt{m}
=\frac{1}{2\sqrt{m}}.
\]

(c):
\\
Let $S = \{x_1, \ldots, x_m\}$ be a sample containing $k$ distinct points $z_1, \ldots, z_k$, where point $z_j$ appears $n_j$ times. We have $n^* = \max_{j=1,\ldots,k} n_j$ and $\sum_{j=1}^k n_j = m$.

For any function $e_z \in \mathcal{E}$ where $e_z(x) = \mathbf{1}_{\{x=z\}}$, we have:
\[
\sum_{i=1}^m \sigma_i e_z(x_i) = 
\begin{cases}
\sum_{i: x_i = z_j} \sigma_i & \text{if } z = z_j \text{ for some } j \in \{1,\ldots,k\} \\
0 & \text{if } z \notin \{z_1,\ldots,z_k\}
\end{cases}
\]

Therefore:
\[
\sup_{e \in \mathcal{E}} \sum_{i=1}^m \sigma_i e(x_i) = \max\left\{0, \max_{j=1,\ldots,k} \sum_{i: x_i = z_j} \sigma_i\right\}
\]

Constructing the finite set, for each $j \in \{1,\ldots,k\}$, define the vector $x^{(j)} \in \mathbb{R}^m$ by:
\[
x^{(j)}_i = \begin{cases} 
1 & \text{if the } i\text{-th sample point equals } z_j \\
0 & \text{otherwise}
\end{cases}
\]

Also define $x^{(0)} = 0 \in \mathbb{R}^m$ (the zero vector).

Let $A = \{x^{(0)}, x^{(1)}, \ldots, x^{(k)}\} \subset \mathbb{R}^m$. Then:
\[
\sup_{e \in \mathcal{E}} \sum_{i=1}^m \sigma_i e(x_i) = \sup_{x \in A} \sum_{i=1}^m \sigma_i x_i
\]

Computing the parameters for Massart's Lemma:

\begin{itemize}
\item The set $A$ is finite with $|A| = k + 1 \leq m + 1$
\item For $j \in \{1,\ldots,k\}$: $\|x^{(j)}\|_2 = \sqrt{n_j} \leq \sqrt{n^*}$
\item For $j = 0$: $\|x^{(0)}\|_2 = 0$
\item Therefore: $R = \max_{x \in A} \|x\|_2 = \max_{j=1,\ldots,k} \sqrt{n_j} = \sqrt{n^*}$
\end{itemize}

By Massart's Lemma, for the finite set $A \subseteq \mathbb{R}^m$ with $R = \max_{x \in A} \|x\|_2 = \sqrt{n^*}$:
\[
\mathbb{E}_{\sigma}\left[\frac{1}{m} \sup_{x \in A} \sum_{i=1}^m \sigma_i x_i\right] \leq \frac{R\sqrt{2\log|A|}}{m} = \frac{\sqrt{n^*} \cdot \sqrt{2\log(k+1)}}{m}
\]

Since $k \leq m$, we have $\log(k+1) \leq \log(m+1)$. Therefore:
\[
\widehat{\mathfrak{R}}_S(\mathcal{E}) = \mathbb{E}_{\sigma}\left[\frac{1}{m} \sup_{e \in \mathcal{E}} \sum_{i=1}^m \sigma_i e(x_i)\right] \leq \frac{\sqrt{2n^* \log(m+1)}}{m}
\]


\end{document}